% ============================================================
%  Chapter 10 — Basis
% ============================================================
\section{Basis}
\label{sec:basis}

A \emph{basis} combines two requirements that are individually familiar: the
vectors span the whole space (so nothing is missing) and they are linearly
independent (so there is no redundancy).  Bases are the workhorses of linear
algebra: they provide coordinates, enable matrix representations, and measure
the size of a space.

% ---- Definition --------------------------------------------
\begin{defbox}[Definition --- Basis of a Vector Space]
\label{def:basis}
A \textbf{basis} for a vector space $V$ over $\F$ is a linearly independent
subset $B\subseteq V$ such that $\spn(B)=V$.
\end{defbox}

% ---- Uniqueness of representation --------------------------
\begin{thmbox}[Theorem --- Unique Representation Characterisation]
\label{thm:basis-unique-rep}
A finite set $\{v_1,\ldots,v_n\}\subseteq V$ is a basis for $V$ if and only if
every vector $v\in V$ can be written \emph{uniquely} as
\[
    v = \alpha_1 v_1 + \alpha_2 v_2 + \cdots + \alpha_n v_n,
    \qquad \alpha_i\in\F.
\]
\end{thmbox}

\begin{proofbox}[Proof --- Unique Representation Characterisation]
($\Rightarrow$) Since the set spans $V$, every $v$ has at least one such
representation.  Suppose $v=\sum\alpha_i v_i=\sum\beta_i v_i$; subtracting
gives $\sum(\alpha_i-\beta_i)v_i=\mathbf{0}$.  Linear independence forces
$\alpha_i=\beta_i$ for all $i$.

($\Leftarrow$) Existence of representations gives $\spn\{v_1,\ldots,v_n\}=V$.
If $\sum\alpha_i v_i=\mathbf{0}$, uniqueness applied to $\mathbf{0}$ shows
$\alpha_i=0$ for all $i$, so the set is l.i.
\hfill$\blacksquare$
\end{proofbox}

% ---- Examples ----------------------------------------------
\begin{exbox}[Example --- Bases in $\R^2$]
\begin{enumerate}
    \item $\{(1,0),(0,1)\}$ is the \emph{standard basis} for $\R^2$.

    \item $S=\{(1,1),(1,-1)\}$: the matrix $A=\begin{pmatrix}1&1\\1&-1\end{pmatrix}$
          has $\rk(A)=2$, so $S$ is l.i.\ and spans $\R^2$.  The decomposition is
          \[
              (x,y) = \tfrac{x+y}{2}(1,1) + \tfrac{x-y}{2}(1,-1).
          \]

    \item $\{(1,0),(0,1),(1,1)\}$ is \emph{not} a basis: it is l.d.\
          since $(1,1)=(1,0)+(0,1)$.
\end{enumerate}
\end{exbox}

% ---- Basis for F^m via rank --------------------------------
\begin{propbox}[Proposition --- Basis for $\F^m$ via Rank]
\label{prop:basis-Fm}
A set $\{v_1,\ldots,v_n\}\subseteq\F^m$ is a basis for $\F^m$ if and only if
$n=m$ and $\rk[v_1\mid\cdots\mid v_n]=n$.
\end{propbox}

\begin{proofbox}[Proof --- Basis for $\F^m$ via Rank]
The set spans $\F^m$ iff $\rk(A)=m$ (Proposition~\ref{prop:rank-generating}),
and is l.i.\ iff $\rk(A)=n$ (Proposition~\ref{prop:rank-li}).
Both conditions together require $m=n$ and $\rk(A)=n$.
\hfill$\blacksquare$
\end{proofbox}

% ---- Characterisation for finite-dimensional spaces --------
\begin{thmbox}[Theorem --- Characterisation of Basis in Finite Dimensions]
\label{thm:basis-char}
Let $V$ be a finite-dimensional vector space with $\dime V=n$, and let
$\{v_1,\ldots,v_n\}\subseteq V$.  The following are equivalent:
\begin{enumerate}[label=(\arabic*)]
    \item $\{v_1,\ldots,v_n\}$ is a basis for $V$.
    \item $V = \spn\{v_1,\ldots,v_n\}$.
    \item $\{v_1,\ldots,v_n\}$ is l.i.
\end{enumerate}
\end{thmbox}

\begin{proofbox}[Proof --- Characterisation of Basis in Finite Dimensions]
(1)$\Rightarrow$(2) and (1)$\Rightarrow$(3) are immediate from the definition.

(2)$\Rightarrow$(1): If $V=\spn\{v_1,\ldots,v_n\}$ and the set were l.d., we
could remove a redundant vector and still span $V$ with fewer than $n$ vectors.
But then any basis of $V$ (which has $n$ elements by definition) would be a
set of $n$ l.i.\ vectors in the span of $n-1$ vectors, contradicting the
Fundamental Lemma (Theorem~\ref{thm:fundamental-lemma}).

(3)$\Rightarrow$(1): An l.i.\ set of $n=\dime V$ vectors in $V$ is
automatically a basis; see the extension theorem
(Theorem~\ref{thm:extend-to-basis}).
\hfill$\blacksquare$
\end{proofbox}

% ---- Extracting a basis from a spanning set ----------------
\begin{thmbox}[Theorem --- Basis Subset of a Spanning Set]
\label{thm:basis-subset}
Any finite spanning set of a vector space $V$ contains a basis for $V$.
More precisely, any finite set of nonzero vectors $\{v_1,\ldots,v_n\}$
contains a basis for $\spn\{v_1,\ldots,v_n\}$.
\end{thmbox}

\begin{propbox}[Proposition --- Finding a Basis via RREF]
\label{prop:basis-via-rref}
Let $A=[v_1\mid\cdots\mid v_n]\in\Mmn{m}{n}$ and let $R$ be its reduced
row-echelon form (RREF).  If $i_1 < i_2 < \cdots < i_r$ are the pivot column
indices of $R$, then $\{v_{i_1},\ldots,v_{i_r}\}$ is a basis for
$\spn\{v_1,\ldots,v_n\}$.
\end{propbox}

\begin{proofbox}[Proof --- Finding a Basis via RREF]
Row operations preserve the column-space relations: a column $v_j$ is a linear
combination of earlier columns iff the corresponding column of $R$ is
non-pivot.  Hence removing non-pivot columns does not change the span.  The
remaining (pivot) columns are l.i.\ because in $R$ they contain the standard
pivot pattern.
\hfill$\blacksquare$
\end{proofbox}

% ---- Extending a l.i. set to a basis -----------------------
\begin{thmbox}[Theorem --- Extending a Linearly Independent Set to a Basis]
\label{thm:extend-to-basis}
In any finite-dimensional vector space $V$ over $\F$, every linearly
independent set $\{v_1,\ldots,v_k\}$ can be extended to a basis of $V$.
That is, there exist $v_{k+1},\ldots,v_n\in V$ such that
$\{v_1,\ldots,v_n\}$ is a basis for $V$.
\end{thmbox}

\begin{proofbox}[Proof --- Extending a Linearly Independent Set to a Basis]
Let $B_0=\{v_1,\ldots,v_k\}$.  If $\spn(B_0)=V$ we are done.  Otherwise
pick $v_{k+1}\in V\setminus\spn(B_0)$; by the sequential criterion
(Proposition~\ref{prop:li-sequential}), $B_1=B_0\cup\{v_{k+1}\}$ is still
l.i.  Repeat.  Since $\dime V$ is finite, this process terminates with a basis.
\hfill$\blacksquare$
\end{proofbox}
